\documentclass[10pt,titlepage]{ltjsarticle}
\title{お酒についてのまとめ}
\author{山田　太郎}
\date{\today}
\usepackage{ascmac}
\usepackage{enumerate}
\begin{document}
\maketitle
\newpage
\tableofcontents
\newpage
\section{日本で飲まれるお酒の種類}
日本では、さまざまなお酒が日々楽しまれいています。その種類には主に以下のようなものがあります。
\begin{screen}
\begin{itemize}
\item
日本酒
\item
焼酎
\begin{itemize}
  \item
  芋焼酎
  \item
  麦焼酎
  \item
  米焼酎 などなど
\end{itemize}
\item
ビール
\item
ワイン
\begin{itemize}
  \item
  赤ワイン
  \item
  白ワイン
\end{itemize}
\item
ウイスキー
\item
カクテル
\item
サワー
\item
その他
\end{itemize}
\end{screen}
このように、多種多様なお酒があります。
\section{醸造酒と蒸留酒}
前節で挙げたお酒は、大きく\textbf{醸造酒}と\textbf{蒸留酒}に分かれます。醸造酒と蒸留酒は、それぞれ以下のように作られるお酒のことです。

\begin{itembox}[l]{醸造酒}
穀類や果物などを発酵させて作る。
\end{itembox}

\begin{itembox}[l]{蒸留酒}
醸造酒を蒸留して作る。
\end{itembox}

前節で挙げたお酒を醸造酒と蒸留酒に分類すると、以下のようになります。
\begin{itemize}
\item
醸造酒
\begin{itemize}
  \item
  日本酒
  \item
  ワイン
  \item
  ビール
\end{itemize}
\item
蒸留酒
\begin{itemize}
  \item
  焼酎
  \item
  ウイスキー
  \item
  ジン、ウォッカなど
\end{itemize}
\end{itemize}
カクテル、サワーについては、焼酎ベースのお酒で作られることが大半なので、基本的には蒸留酒と分類されることが多いです。
\section{甲類焼酎と乙類焼酎}
焼酎はさらに、{\bf 甲類焼酎}と{\bf 乙類焼酎}に分類されます。
\begin{itembox}[l]{甲類焼酎}
連続蒸留によって作られる、アルコール度数36パーセント以下の焼酎。
\end{itembox}

\begin{itembox}[l]{乙類焼酎}
単式蒸留によって作られる、アルコール度数45パーセント以下の焼酎。
\end{itembox}

この分類は、日本における{\bf 酒税法}との関係で生まれたものです。

\section{日本の人気のお酒ランキング}
そんなさまざまなお酒ですが、日本の人気ランキングを発表しましょう。
\begin{itembox}[l]{日本で人気のお酒ランキング}
\begin{enumerate}[第1位.]
\item
ビール
\item
日本酒
\item
ワイン (特に白ワイン)
\end{enumerate}
\end{itembox}
\subsection{海外での人気のお酒}
海外では、国によって人気のお酒が異なりますが、アメリカではやはりビールが最も人気のお酒です。{\bf 「とりあえずビール」}という言葉が象徴するように、やはり世界中で愛されているのですね。
\section{まとめ}
お酒は時に人と人とのコミュニケーションを楽しく、円滑に彩ってくれますが、飲み過ぎによるトラブル、健康被害も多数報告されています。
\begin{center}
\Large{\textgt{お酒は楽しく、ほどほどに！}}
\end{center}
これを肝に銘じて、楽しくお酒と付き合っていきましょう。
\end{document}

